\section{Discussion}

This work demonstrates the feasibility of applying a retrieval and reranking architecture for regulatory document entity linking. By combining a bi-encoder for efficient candidate retrieval with a cross-encoder for fine-grained candidate selection, the proposed framework is able to address the two main challenges of the task: identifying relevant candidates from a large regulatory corpus and distinguishing between semantically similar regulatory entities.

The effectiveness of transformer-based architectures for this task can be further understood through the connection between modern Hopfield networks and attention mechanisms. As discussed in the theoretical foundation of this work, the attention mechanism employed by transformer architectures can be interpreted through the energy formulation of modern Hopfield networks, where attention performs an associative retrieval process over stored representations. From this perspective, regulatory document entity linking can be viewed as an associative retrieval problem, where regulatory references act as queries and regulatory entities represent stored patterns within the learned representation space.

This perspective is particularly relevant for regulatory document entity linking, as references to the same regulatory document entity may exhibit substantial linguistic variation while still representing the same underlying concept. Rather than relying solely on direct lexical similarity, attention-based models are able to learn contextual associations between regulatory references and candidate entities. This enables the system to identify relationships despite differences in wording, document structure, and available contextual information.

However, the quality of these learned associations remains dependent on the diversity and coverage of the underlying training data. As demonstrated in the error analysis, limitations in the available dataset, including heuristic labelling, constrained augmentation, and the focus on EUR-Lex-based regulatory information, may restrict the ability of the model to generalize beyond the observed data distribution. Therefore, improving regulatory document entity linking is not solely dependent on increasing model complexity, but also requires broader and more representative regulatory knowledge coverage.

The two-stage architecture introduces an additional trade-off between retrieval efficiency and linking accuracy. The bi-encoder enables scalable retrieval from a large candidate space, while the cross-encoder provides more precise discrimination between retrieved candidates. However, this separation also introduces distinct failure modes, as errors during candidate retrieval cannot be recovered during reranking. This highlights the importance of balancing retrieval recall with reranking precision when designing entity linking systems for large regulatory corpora.

Overall, the results indicate that attention-based retrieval and reranking architectures provide a suitable foundation for regulatory document entity linking. However, the ability of such systems to generalize effectively depends not only on the underlying model architecture but also on the quality, diversity, and representativeness of the regulatory information used to construct the learned representation space.


\section{Future Work}

Several directions for future work can be identified based on the limitations and observations from this work.

One potential extension is the expansion of the dataset beyond the current EUR-Lex-based regulatory corpus. While EUR-Lex provides a large and structured source of European regulatory information, incorporating additional jurisdictions and regulatory authorities would increase the diversity of legal terminology, document structures, and regulatory relationships represented within the training data. Such an extension would require an analysis of how different jurisdictions describe equivalent regulatory concepts and relationships. Based on this analysis, representative examples from additional jurisdictions could be incorporated into the training dataset to improve the coverage of the learned representation space.

Alternatively, a normalization or adapter layer could be investigated to address differences in jurisdiction-specific terminology. Rather than requiring the model to learn all variations directly from examples, such a layer could transform jurisdiction-specific expressions into a shared regulatory representation, allowing the entity linking model to operate across multiple legal domains while maintaining consistent behaviour.

Another direction for future work is the continuous refinement of the linking model through hard negative mining. During inference, incorrect high-scoring candidates provide valuable information about difficult distinctions within the candidate space. These hard negatives could be collected and incorporated into subsequent training iterations, allowing the model to better differentiate between highly similar regulatory entities. This approach could improve the decision boundaries of the cross-encoder and refine the associative representations learned by the model.

Additionally, the linked regulatory entities produced by the system provide an opportunity to develop higher-level knowledge representation systems. Once regulatory references are reliably linked to their corresponding entities, these relationships can be used to construct or enhance regulatory knowledge graphs. Such graph-based representations could enable further applications, including regulatory dependency analysis, document relationship traversal, impact analysis of regulatory changes, and exploration of interconnected legal frameworks.

Finally, future work could investigate more dynamic approaches for maintaining regulatory document entity representations as new documents and relationships become available. Regulatory information continuously evolves through amendments, replacements, and newly introduced legislation. Developing methods for incorporating newly observed regulatory information while preserving existing knowledge could allow the system to adapt to changing regulatory environments without requiring complete retraining.